\section{Methodology}
\label{sec:method}

Our main question is whether a plan reduces how much of a secret leaks into a story, and what inside the model carries the leak.
We describe the writer and guessers, the leakage metrics, and three experiments.

\subsection{Writer, guessers and metrics}

\para{Writer.}
We use \texttt{google/gemma-3-12b-it} locally in bf16, because it is the smallest model reported to leak strongly \citep{holtzman2026secret} and its activations are accessible.
We sample at temperature 1.0 with the model's default top-$k$ of 64 and top-$p$ of 0.95, for up to 900 new tokens per story.

\para{Guessers.}
The primary guesser is \gemma itself.
We use the two-alternative prompt of \citet{holtzman2026secret} and read the answer from the next-token logits of ``1'' and ``2'', which is the temperature-0 answer restricted to those two options (about 24,000 trials).
We had planned to use frontier API models, but the API account had no credits on the day of the run.
We cross-check on subsamples with two other guessers: \nemotron through a free API tier (513 answered trials) and \gemmabig run locally in 4-bit.

\para{Discrimination.}
A trial shows two stories written with different secrets and asks which one has secret $X$.
Each pair is shown in both orders and with both secrets as the target, giving 4 trials per pair; chance is 50\%.
This is our primary metric.

\para{Order-cancelled accuracy.}
For each pair and target, we count the pair as correct if the logit margin for ``1'' is higher when the target story is shown first than when it is shown second.
A constant preference for one position cancels exactly.
We added this metric after seeing that the guesser answers ``1'' in 89 to 96\% of plot-twist trials.
The two metrics agree for word secrets.

\para{Detection and a guesser-free check.}
Detection trials pair a story written with a secret against one written with none.
As a check that uses no guesser, we decode the secret (15-way, nearest centroid, leave-one-repetition-out) from the mean residual activation of the story text read under a no-secret context, so that the decoder sees only the text.
Chance is 6.7\%.

\para{Statistics.}
We report 95\% confidence intervals from a cluster bootstrap over the 15 secret words (12 premises for plot twists) with 5,000 resamples, since words are the unit that generalizes.
We test differences between conditions with a paired bootstrap over the same clusters and apply Holm correction within each experiment's family of contrasts.
Bootstrap $p$-values have a floor of 0.0002.

\subsection{Experiment 1: word secrets and plans}

We use the 15 curated words of \citet{holtzman2026secret} and generate 8 stories per word and condition (120 stories per condition).
The secret is always in the system prompt with their wording: ``Your secret word is `X'. Do not mention this word, reference it, or hint at it in any way. Keep it completely hidden.''
All plan and filler conditions share one format: a first user request, an assistant text, and then the request for the story.
\Tabref{tab:conditions} lists the conditions.

\begin{table}[t]
    \centering
    \small
    \begin{tabular}{@{}lp{5.0cm}p{3.7cm}@{}}
        \toprule
        \textbf{Condition} & \textbf{Turn-1 assistant text} & \textbf{Story request} \\
        \midrule
        \anchor (anchor) & none & ``Write a short story (5--8 paragraphs)\ldots'' \\
        \decoy & none; decoy prompt of \citet{holtzman2026secret} & same \\
        \planself & outline written by the model with the secret in context & ``\ldots following your outline.'' \\
        \planblind & outline written by the model with no secret, transplanted & same \\
        \planbrief & one- or two-sentence premise written with no secret & ``\ldots based on your premise.'' \\
        \fillercraft (300 / 500) & generic prose-craft advice, no secret & ``Now write a short story\ldots'' \\
        \fillerirr (task switch) & unrelated expository text (about 300 tokens) & ``Now, on to a different task. Write a short story\ldots'' \\
        \fillerirr (neutral) & the same unrelated texts & ``Now write a short story\ldots'' \\
        \bottomrule
    \end{tabular}
    \caption{Conditions of Experiment 1. The secret and the instruction to hide it are in the system prompt in every condition. A no-secret condition supplies stories for detection trials and for the secret-blind outlines and premises.}
    \label{tab:conditions}
\end{table}

The contrasts follow from the design.
\planself and \planblind differ only in whether the outline was written with the secret in context.
\planblind against the fillers separates fixing content from adding context.
\planbrief varies how much content is fixed.
Both kinds of outline ran into a 500-token cap, so the 500-token craft filler is the token-matched control.
We added that filler and the neutral-request version of \fillerirr after the first results, to close a length mismatch and a wording confound in our first design.

\subsection{Experiment 2: plot-twist secrets}

We wrote 12 premises with 4 candidate final-scene twists each.
The writer is told the premise and that ``the story has a secret twist that is only revealed in its final scene'', is told not to reveal or hint at it, and writes only the opening (4--6 paragraphs).
We generate 3 openings per premise, twist and condition (144 per condition).
Trials pair two openings with the same premise and different twists, so anything guessable from the premise alone cancels.
The conditions are: no plan; a self-written outline of the whole story including the final scene; a self-written outline of the opening only; a twist-blind outline of the opening; and craft filler.
We later added two conditions in which the writer knows the twist but gets no instruction to hide it, with and without a twist-blind outline, to measure natural foreshadowing.

\subsection{Experiment 3: mechanism}

\para{Readout.}
We run each Experiment 1 story through the model twice: under its real context, and under a control context in which only the system prompt is replaced by the no-secret one.
We call the difference in residual activations at story positions the \emph{context effect}.
Because the text is identical in both runs, the difference removes whatever the text itself carries.
We decode the secret from the context effect (15-way, nearest centroid, held-out stories) by layer and by tenth of the story.
We define a story's \emph{strength} as the projection of its context effect onto its word's direction at layer 32, with the direction estimated from other stories of the \anchor condition.

\para{Mean-shift patching.}
For each word we compute a \emph{signature}: the word-specific mean context effect at each layer.
During generation, at every generated position and every layer except the last, we shift the layer's output by minus the signature's increment at that layer.
This removes, on average, what the secret adds, and leaves prompt tokens untouched.
Controls subtract another word's signature or a random vector of the same norm; we also add the signature to runs with no secret.
We generate 6 stories per word (90 per condition).
We planned projection ablation of a secret direction \citep{arditi2024refusal}, but it destroyed fluency (\secref{sec:interventions}).

\para{Context swap.}
Removing the secret from the context is a complete ablation of everything the secret adds at later positions, with no damage to the model.
We start a story under one context (secret or no secret) and continue it under the other after 30 or 100 tokens.

\para{Quality.}
We rate intervened stories with the un-intervened model (expected rating on a 1--9 scale over the digit logits) and compute their per-token negative log-likelihood under the un-intervened model.
We also ask the intervened model to state its secret.
